% ==============================================================================
% RÉSUMÉ (FRANÇAIS)
% ==============================================================================
\chapter*{Résumé}
\addcontentsline{toc}{chapter}{Résumé}

\begin{tcolorbox}[
  enhanced,
  colback=bleuprincipal!5!white,
  colframe=bleuprincipal,
  boxrule=1.5pt, arc=6pt,
  left=10pt, right=10pt, top=8pt, bottom=8pt
]

Ce mémoire présente la conception et le développement d'un système complet
d'analyse automatique de sentiments appliqué aux commentaires d'étudiants
sur leurs cours universitaires. Face à la difficulté de traiter manuellement
des milliers de retours étudiants, ce projet propose une solution intelligente,
automatisée et temps réel, exploitant les avancées récentes du Traitement
Automatique du Langage Naturel (TALN).

\vspace{0.5em}

Le système repose sur le modèle de langage \textbf{CamemBERT}
(\textit{camembert-base}), une variante française de BERT développée par
Inria et CNRS, fine-tuné pour la classification de sentiments en trois
catégories : \textbf{Négatif}, \textbf{Neutre} et \textbf{Positif}. Le modèle
a été entraîné sur un jeu de données de 225 commentaires étudiants en langue
française, équilibrés entre les trois classes, avec les hyperparamètres
optimaux suivants : taux d'apprentissage de $2 \times 10^{-5}$, taille de
lot de 16, 3 epochs d'entraînement et une longueur maximale de séquence de
128 tokens. L'optimiseur AdamW avec un scheduler linéaire à palier de
chauffe a été utilisé pour garantir la convergence. Le modèle atteint
une précision de \textbf{79,41\%} sur le jeu de test.

\vspace{0.5em}

L'architecture globale du système comprend quatre composants principaux :
(1) un \textbf{backend RESTful} développé avec FastAPI et SQLAlchemy
communiquant avec une base de données PostgreSQL pour la persistance des
analyses ; (2) un \textbf{tableau de bord analytique} Streamlit enrichi
de visualisations interactives Plotly pour le suivi en temps réel des
tendances ; (3) une \textbf{application mobile} Flutter offrant aux
étudiants une interface moderne et intuitive pour soumettre leurs commentaires
et consulter les résultats ; (4) un \textbf{pipeline de prétraitement}
NLP complet comprenant la normalisation Unicode, la suppression des caractères
spéciaux et la tokenisation CamemBERT.

\vspace{0.5em}

Les résultats obtenus démontrent la faisabilité et l'efficacité d'une telle
approche, même avec un jeu de données limité, ouvrant la voie à des
améliorations significatives avec davantage de données annotées.

\end{tcolorbox}

\vspace{0.8em}
\noindent\textbf{Mots-clés :} Analyse de sentiments, TALN, CamemBERT,
Apprentissage automatique, FastAPI, Flutter, PostgreSQL, Classification de
texte, Fine-tuning, Deep learning.

\cleardoublepage
